% Fuente: Auxiliar 6, Problema 1
\begingroup
\color{MidnightBlue}
\begin{enumerate}
\item
Los vértices del conjunto factible son:
\begin{equation}
O=(0,0),\qquad
B=(0,3),\qquad
E=\left(\frac{10}{3},\frac{4}{3}\right),
\qquad
C=(4,0).
\end{equation}

\begin{figure}
    \centering
    \includegraphics[width=0.5\linewidth]{ejercicios/Capitulo-4/ejercicio_4_1_conjunto_factible.png}
    \caption{Conjunto Factible P1}
    \label{ex:4:1:fig:placeholder}
\end{figure}

Agregando variables de holgura $s_1,s_2\geq 0$ y pasando a minimización, se obtiene la forma estándar:
\begin{equation}
\begin{aligned}
\min \quad -2x_1-3x_2 \\
x_1+2x_2+s_1 &= 6,\\
2x_1+x_2+s_2 &= 8,\\
x_1,x_2,s_1,s_2 &\geq 0.
\end{aligned}
\end{equation}

Así,
\begin{equation}
A=
\begin{pmatrix}
1 & 2 & 1 & 0\\
2 & 1 & 0 & 1
\end{pmatrix},
\qquad
b=
\begin{pmatrix}
6\\
8
\end{pmatrix},
\qquad
c=
\begin{pmatrix}
-2\\
-3\\
0\\
0
\end{pmatrix}.
\end{equation}

Las soluciones básicas se obtienen escogiendo $m=2$ variables básicas entre las $4$ variables
$\{x_1,x_2,s_1,s_2\}$.

\begin{center}
\begin{tabular}{c|c|c|c}
Base & Variables no básicas & Solución $(x_1,x_2,s_1,s_2)$ & ¿Factible?\\
\hline
$\{s_1,s_2\}$ & $x_1=x_2=0$ & $(0,0,6,8)$ & Sí\\
$\{x_2,s_2\}$ & $x_1=s_1=0$ & $(0,3,0,5)$ & Sí\\
$\{x_1,s_1\}$ & $x_2=s_2=0$ & $(4,0,2,0)$ & Sí\\
$\{x_1,x_2\}$ & $s_1=s_2=0$ & $\left(\frac{10}{3},\frac{4}{3},0,0\right)$ & Sí\\
$\{x_1,s_2\}$ & $x_2=s_1=0$ & $(6,0,0,-4)$ & No\\
$\{x_2,s_1\}$ & $x_1=s_2=0$ & $(0,8,-10,0)$ & No
\end{tabular}
\end{center}

Por lo tanto, las soluciones básicas factibles son:
\begin{equation}
(0,0,6,8),\qquad
(0,3,0,5),\qquad
(4,0,2,0),\qquad
\left(\frac{10}{3},\frac{4}{3},0,0\right).
\end{equation}

Evaluando la función objetivo original:
\begin{equation}
\begin{array}{c|c}
(x_1,x_2) & z=2x_1+3x_2\\
\hline
(0,0) & 0\\
(0,3) & 9\\
(4,0) & 8\\
\left(\frac{10}{3},\frac{4}{3}\right) & \frac{32}{3}
\end{array}
\end{equation}

Luego, la solución óptima es:
\begin{equation}
x^*=\left(\frac{10}{3},\frac{4}{3}\right),
\qquad
z^*=\frac{32}{3}.
\end{equation}

\item
\textbf{Iteración 0}

Base inicial:
\begin{equation}
B_0=\{s_1,s_2\},
\qquad
x_B=
\begin{pmatrix}
s_1\\
s_2
\end{pmatrix}
=
\begin{pmatrix}
6\\
8
\end{pmatrix}.
\end{equation}

La solución inicial es:
\begin{equation}
(x_1,x_2,s_1,s_2)=(0,0,6,8),
\end{equation}
que corresponde al punto
\begin{equation}
O=(0,0).
\end{equation}

Como $c_B=(0,0)$, se tiene:
\begin{equation}
p^T=c_B^T B^{-1}=(0,0).
\end{equation}

Los costos reducidos son:
\begin{equation}
\bar{c}_{x_1}=c_{x_1}-p^TA_1=-2,
\qquad
\bar{c}_{x_2}=c_{x_2}-p^TA_2=-3.
\end{equation}

 En este caso, buscaremos la direccion de mayor desplazamiento. Como $\bar{c}_{x_2}$ es el más negativo, entra $x_2$ a la base.

La dirección es:
\begin{equation}
d_B=-B^{-1}A_2
=
-
\begin{pmatrix}
1 & 0\\
0 & 1
\end{pmatrix}
\begin{pmatrix}
2\\
1
\end{pmatrix}
=
\begin{pmatrix}
-2\\
-1
\end{pmatrix}.
\end{equation}

Luego:
\begin{equation}
d=(0,1,-2,-1).
\end{equation}

El máximo paso factible es:
\begin{equation}
\theta^*
=
\min\left\{
\frac{6}{2},
\frac{8}{1}
\right\}
=
3.
\end{equation}

Por lo tanto, sale $s_1$ de la base. La nueva solución es:
\begin{equation}
x^{(1)}
=
(0,3,0,5),
\end{equation}
que corresponde al punto:
\begin{equation}
B=(0,3).
\end{equation}

\textbf{Iteración 1}

La nueva base es:
\begin{equation}
B_1=\{x_2,s_2\},
\qquad
B=
\begin{pmatrix}
2 & 0\\
1 & 1
\end{pmatrix},
\qquad
B^{-1}
=
\begin{pmatrix}
\frac{1}{2} & 0\\
-\frac{1}{2} & 1
\end{pmatrix}.
\end{equation}

Además,
\begin{equation}
x_B=
\begin{pmatrix}
x_2\\
s_2
\end{pmatrix}
=
\begin{pmatrix}
3\\
5
\end{pmatrix},
\qquad
c_B=
\begin{pmatrix}
-3\\
0
\end{pmatrix}.
\end{equation}

Entonces:
\begin{equation}
p^T=c_B^T B^{-1}
=
\begin{pmatrix}
-3 & 0
\end{pmatrix}
\begin{pmatrix}
\frac{1}{2} & 0\\
-\frac{1}{2} & 1
\end{pmatrix}
=
\begin{pmatrix}
-\frac{3}{2} & 0
\end{pmatrix}.
\end{equation}

Los costos reducidos de las variables no básicas $x_1$ y $s_1$ son:
\begin{equation}
\bar{c}_{x_1}
=
-2-
\begin{pmatrix}
-\frac{3}{2} & 0
\end{pmatrix}
\begin{pmatrix}
1\\
2
\end{pmatrix}
=
-\frac{1}{2},
\end{equation}

\begin{equation}
\bar{c}_{s_1}
=
0-
\begin{pmatrix}
-\frac{3}{2} & 0
\end{pmatrix}
\begin{pmatrix}
1\\
0
\end{pmatrix}
=
\frac{3}{2}.
\end{equation}

Como $\bar{c}_{x_1}<0$, entra $x_1$ a la base.

La dirección básica se calcula como:
\begin{equation}
B^{-1}A_1
=
\begin{pmatrix}
\frac{1}{2} & 0\\
-\frac{1}{2} & 1
\end{pmatrix}
\begin{pmatrix}
1\\
2
\end{pmatrix}
=
\begin{pmatrix}
\frac{1}{2}\\
\frac{3}{2}
\end{pmatrix}.
\end{equation}

El paso máximo factible es:
\begin{equation}
\theta^*
=
\min\left\{
\frac{3}{1/2},
\frac{5}{3/2}
\right\}
=
\min\left\{
6,\frac{10}{3}
\right\}
=
\frac{10}{3}.
\end{equation}

Por lo tanto, sale $s_2$ de la base. La nueva solución es:
\begin{equation}
x^{(2)}
=
\left(
\frac{10}{3},
\frac{4}{3},
0,
0
\right),
\end{equation}
que corresponde al punto:
\begin{equation}
E=\left(\frac{10}{3},\frac{4}{3}\right).
\end{equation}

\textbf{Iteración 2}

La base óptima es:
\begin{equation}
B_2=\{x_2,x_1\}.
\end{equation}

Se tiene:
\begin{equation}
B=
\begin{pmatrix}
2 & 1\\
1 & 2
\end{pmatrix},
\qquad
B^{-1}
=
\begin{pmatrix}
\frac{2}{3} & -\frac{1}{3}\\
-\frac{1}{3} & \frac{2}{3}
\end{pmatrix}.
\end{equation}

Además, como el orden de la base es $\{x_2,x_1\}$, se tiene:
\begin{equation}
c_B=
\begin{pmatrix}
-3\\
-2
\end{pmatrix}.
\end{equation}

Luego:
\begin{equation}
p^T=c_B^T B^{-1}
=
\begin{pmatrix}
-3 & -2
\end{pmatrix}
\begin{pmatrix}
\frac{2}{3} & -\frac{1}{3}\\
-\frac{1}{3} & \frac{2}{3}
\end{pmatrix}
=
\begin{pmatrix}
-\frac{4}{3} & -\frac{1}{3}
\end{pmatrix}.
\end{equation}

Los costos reducidos de las variables no básicas $s_1$ y $s_2$ son:
\begin{equation}
\bar{c}_{s_1}
=
0-
\begin{pmatrix}
-\frac{4}{3} & -\frac{1}{3}
\end{pmatrix}
\begin{pmatrix}
1\\
0
\end{pmatrix}
=
\frac{4}{3}>0,
\end{equation}

\begin{equation}
\bar{c}_{s_2}
=
0-
\begin{pmatrix}
-\frac{4}{3} & -\frac{1}{3}
\end{pmatrix}
\begin{pmatrix}
0\\
1
\end{pmatrix}
=
\frac{1}{3}>0.
\end{equation}

Como todos los costos reducidos son no negativos, la solución es óptima para el problema de minimización equivalente. Por lo tanto, para el problema original de maximización:
\begin{equation}
x^*
=
\left(
\frac{10}{3},
\frac{4}{3}
\right),
\qquad
z^*
=
2\cdot \frac{10}{3}
+
3\cdot \frac{4}{3}
=
\frac{32}{3}.
\end{equation}

La trayectoria del método Simplex fue:
\begin{equation}
(0,0)
\longrightarrow
(0,3)
\longrightarrow
\left(\frac{10}{3},\frac{4}{3}\right).
\end{equation}

En el método Simplex Revisado no se recalcula toda la base desde cero,
sino que se actualiza $B^{-1}$ usando la columna entrante. La idea es formar
la matriz aumentada
\begin{equation}
\left[B^{-1}\mid u\right],
\qquad
u=B^{-1}A_j,
\end{equation}
donde $A_j$ es la columna de la variable que entra a la base.

Luego, mediante operaciones elementales por filas, se transforma $u$ en el
vector canónico $e_l$, donde $l$ es la posición en el vector base de la variable que sale de la base. El bloque izquierdo
resultante corresponde a la nueva matriz $B^{-1}$.

\textbf{Iteración 0}

La base inicial es
\begin{equation}
B_0=\{s_1,s_2\}.
\end{equation}

Por lo tanto,
\begin{equation}
B_0=
\begin{pmatrix}
1 & 0\\
0 & 1
\end{pmatrix},
\qquad
B_0^{-1}=
\begin{pmatrix}
1 & 0\\
0 & 1
\end{pmatrix}.
\end{equation}

En esta iteración entra $x_2$, cuya columna es
\begin{equation}
A_2=
\begin{pmatrix}
2\\
1
\end{pmatrix}.
\end{equation}

Entonces,
\begin{equation}
u=B_0^{-1}A_2
=
\begin{pmatrix}
1 & 0\\
0 & 1
\end{pmatrix}
\begin{pmatrix}
2\\
1
\end{pmatrix}
=
\begin{pmatrix}
2\\
1
\end{pmatrix}.
\end{equation}

Como el test de razón mínima indica que sale $s_1$, la fila pivote es la
primera. Por lo tanto, se debe transformar $u$ en
\begin{equation}
e_1=
\begin{pmatrix}
1\\
0
\end{pmatrix}.
\end{equation}

Formamos la matriz aumentada:
\begin{equation}
\left[
B_0^{-1}\mid u
\right]
=
\left[
\begin{array}{cc|c}
1 & 0 & 2\\
0 & 1 & 1
\end{array}
\right].
\end{equation}

Aplicamos operaciones elementales para convertir la última columna en $e_1$.

Primero eliminamos el elemento de la segunda fila:
\begin{equation}
F_2 \leftarrow F_2-\frac{1}{2}F_1.
\end{equation}

Así,
\begin{equation}
\left[
\begin{array}{cc|c}
1 & 0 & 2\\
0 & 1 & 1
\end{array}
\right]
\longrightarrow
\left[
\begin{array}{cc|c}
1 & 0 & 2\\
-\frac{1}{2} & 1 & 0
\end{array}
\right].
\end{equation}

Luego normalizamos la fila pivote:
\begin{equation}
F_1 \leftarrow \frac{1}{2}F_1.
\end{equation}

Entonces,
\begin{equation}
\left[
\begin{array}{cc|c}
1 & 0 & 2\\
-\frac{1}{2} & 1 & 0
\end{array}
\right]
\longrightarrow
\left[
\begin{array}{cc|c}
\frac{1}{2} & 0 & 1\\
-\frac{1}{2} & 1 & 0
\end{array}
\right].
\end{equation}

Como la última columna quedó igual a $e_1$, el bloque izquierdo es la nueva
inversa de la base:
\begin{equation}
B_1^{-1}
=
\begin{pmatrix}
\frac{1}{2} & 0\\
-\frac{1}{2} & 1
\end{pmatrix}.
\end{equation}

La nueva base es
\begin{equation}
B_1=\{x_2,s_2\}.
\end{equation}

\textbf{Iteración 1}

Ahora entra $x_1$, cuya columna es
\begin{equation}
A_1=
\begin{pmatrix}
1\\
2
\end{pmatrix}.
\end{equation}

Calculamos:
\begin{equation}
u=B_1^{-1}A_1
=
\begin{pmatrix}
\frac{1}{2} & 0\\
-\frac{1}{2} & 1
\end{pmatrix}
\begin{pmatrix}
1\\
2
\end{pmatrix}
=
\begin{pmatrix}
\frac{1}{2}\\
\frac{3}{2}
\end{pmatrix}.
\end{equation}

El test de razón mínima indica que sale $s_2$, es decir, la fila pivote es la
segunda. Por lo tanto, ahora se debe transformar $u$ en
\begin{equation}
e_2=
\begin{pmatrix}
0\\
1
\end{pmatrix}.
\end{equation}

Formamos la matriz aumentada:
\begin{equation}
\left[
B_1^{-1}\mid u
\right]
=
\left[
\begin{array}{cc|c}
\frac{1}{2} & 0 & \frac{1}{2}\\
-\frac{1}{2} & 1 & \frac{3}{2}
\end{array}
\right].
\end{equation}

Primero normalizamos la fila pivote:
\begin{equation}
F_2 \leftarrow \frac{2}{3}F_2.
\end{equation}

Así,
\begin{equation}
\left[
\begin{array}{cc|c}
\frac{1}{2} & 0 & \frac{1}{2}\\
-\frac{1}{2} & 1 & \frac{3}{2}
\end{array}
\right]
\longrightarrow
\left[
\begin{array}{cc|c}
\frac{1}{2} & 0 & \frac{1}{2}\\
-\frac{1}{3} & \frac{2}{3} & 1
\end{array}
\right].
\end{equation}

Luego eliminamos el elemento de la primera fila:
\begin{equation}
F_1 \leftarrow F_1-\frac{1}{2}F_2.
\end{equation}

Por lo tanto,
\begin{equation}
\left[
\begin{array}{cc|c}
\frac{1}{2} & 0 & \frac{1}{2}\\
-\frac{1}{3} & \frac{2}{3} & 1
\end{array}
\right]
\longrightarrow
\left[
\begin{array}{cc|c}
\frac{2}{3} & -\frac{1}{3} & 0\\
-\frac{1}{3} & \frac{2}{3} & 1
\end{array}
\right].
\end{equation}

Como la última columna quedó igual a $e_2$, el bloque izquierdo corresponde a
la nueva inversa de la base:
\begin{equation}
B_2^{-1}
=
\begin{pmatrix}
\frac{2}{3} & -\frac{1}{3}\\
-\frac{1}{3} & \frac{2}{3}
\end{pmatrix}.
\end{equation}

La nueva base es
\begin{equation}
B_2=\{x_2,x_1\}.
\end{equation}

\textbf{Iteración 2}
Dado que los costos reducidos de las variables no básicas $s_1$ y $s_2$ son positivas, la solución encontrada es óptima para el problema.

\item
El primal está en la forma:
\begin{equation}
\begin{aligned}
\max \quad & c^Tx\\
\text{s.a.}\quad & Ax\leq b,\\
& x\geq 0.
\end{aligned}
\end{equation}

Por lo tanto, el dual es:
\begin{equation}
\begin{aligned}
\min \quad & b^Ty\\
\text{s.a.}\quad & A^Ty\geq c,\\
& y\geq 0.
\end{aligned}
\end{equation}

En este caso:
\begin{equation}
\begin{aligned}
\min \quad & 6y_1+8y_2\\
\text{s.a.}\quad
& y_1+2y_2 \geq 2,\\
& 2y_1+y_2 \geq 3,\\
& y_1,y_2 \geq 0.
\end{aligned}
\end{equation}

Por Dualidad Fuerte, el óptimo dual de coincidir con el óptimo de modo que:
\begin{equation}
    z^*=\frac{32}{3}=b^Ty^* 
\end{equation}

Durante el auxiliar grabado, me falto mostrar que es posible obtener los valores de $y^*$ mediante las condiciones de holgura complementaria:
\begin{equation}
\begin{aligned}
y_1(6-x_1-2x_2) &= 0,\\
y_2(8-2x_1-x_2) &= 0,\\
x_1(y_1+2y_2-2) &= 0,\\
x_2(2y_1+y_2-3) &= 0.
\end{aligned}
\end{equation}

Como en el óptimo primal se tiene:
\begin{equation}
x_1^*>0,
\qquad
x_2^*>0,
\end{equation}
por holgura complementaria las restricciones duales asociadas a $x_1$ y $x_2$ son activas:
\begin{equation}
\begin{aligned}
y_1+2y_2 &= 2,\\
2y_1+y_2 &= 3.
\end{aligned}
\end{equation}

Resolviendo:
\begin{equation}
y_1=\frac{4}{3},
\qquad
y_2=\frac{1}{3}.
\end{equation}

El valor óptimo dual es:
\begin{equation}
w^*
=
6\cdot \frac{4}{3}
+
8\cdot \frac{1}{3}
=
8+\frac{8}{3}
=
\frac{32}{3}.
\end{equation}
\end{enumerate}
\endgroup
